\section{Описание}
Требуется реализовать словарь на основе сбалансированного АВЛ-дерева поиска. 

Как сказано в \cite{Kormen}, деревья поиска позволяют выполнять операции вставки, удаления и поиска за время, пропорциональное высоте дерева. АВЛ-дерево строго контролирует свою высоту: для каждого узла разность высот его левого и правого поддеревьев (баланс) не превышает 1 по модулю. При нарушении этого свойства выполняются малые или большие вращения (rotations). Это гарантирует логарифмическую асимптотику $O(\log N)$ для основных операций в худшем случае.

В данной работе алгоритм реализован с использованием современных стандартов C++:
\begin{enumerate}
\item Узлы дерева управляются через умные указатели \texttt{std::unique\_ptr}, что гарантирует отсутствие утечек памяти при удалениях узлов или перехвате исключений \texttt{std::bad\_alloc}.
\item Нормализация ключей (приведение к нижнему регистру) реализована in-place с помощью быстрой ASCII-арифметики, избегая выделения памяти под новые строки.
\item Для реализации команд \texttt{Save} и \texttt{Load} используется бинарный формат. Дерево сохраняется методом прямого обхода (Pre-order). При загрузке это позволяет восстановить структуру дерева за $O(N)$ без необходимости выполнять балансировку. 
\item При загрузке словаря данные сначала помещаются во временный корень. Текущий словарь подменяется новым только в том случае, если чтение файла прошло без ошибок, что удовлетворяет строгим требованиям технического задания.
\end{enumerate}

\section{Исходный код}
Программа состоит из класса \texttt{AVLTree} и функции \texttt{main}, реализующей парсинг команд.

\textbf{Структура узла дерева:}
\begin{lstlisting}[language=C++]
struct TNode {
    std::string key;
    uint64_t value;
    uint8_t height;
    std::unique_ptr<TNode> left;
    std::unique_ptr<TNode> right;

    TNode(std::string k, uint64_t v) 
        : key(std::move(k)), value(v), height(1), 
          left(nullptr), right(nullptr) {}
};
\end{lstlisting}

Вращения реализованы через функцию \texttt{std::move}, которая передает владение узлами между умными указателями без реального копирования данных в памяти.

Полный текст программы приведен в приложении к отчету (или в листинге ниже):

\lstinputlisting[language=C++, caption={Реализация словаря на АВЛ-дереве}, label={lst:maincpp}]{../recursive_avl.cpp}
\pagebreak

\section{Консоль}
В данном примере приведена компиляция и запуск программы, а также проверка корректности вывода с помощью утилиты \texttt{diff}. Отсутствие вывода после вызова команды \texttt{diff} означает полное (побайтовое) совпадение ожидаемого и реального результатов.

\begin{alltt}
$ g++ -std=c++17 -Wall -O3 recursive_avl.cpp -o avl_dict
$ cat test.in
+ Apple 100
+ banana 200
+ APPLE 400
apple
- banana
banana
! Save my_dict.bin
- apple
! Load my_dict.bin
apple
$ ./avl_dict < test.in > result.out
$ cat expected.out
OK
OK
Exist
OK: 100
OK
NoSuchWord
OK
OK
OK
OK: 100
$ diff expected.out result.out
$ 
\end{alltt}
\pagebreak
